% Vietnamese translation for OI Public Library
% Source-PDF: IOI中国国家候选队论文/2008/Day1/7.方戈《浅析信息学竞赛中一类与物理有关的问题》/water tanks.pdf
\documentclass[11pt]{article}
\usepackage{oipl-vn}

\title{Bể nước}
\author{ACM ICPC World Finals Tokyo 2006--2007}
\date{}

\begin{document}
\maketitle

\origtitle{3810 - Problem: Water Tanks}
\sourcepath{ioi/2008/water-tanks.pdf}

\section*{Đề bài}

Công ty Aqua Container Management quản lý các cơ sở chứa nước. Họ đang cân
nhắc một hệ thống lưu trữ nước gồm một dãy các bể thẳng đứng được nối với
nhau. Mỗi bể có diện tích mặt cắt ngang nằm ngang bằng một mét vuông, nhưng
các bể có chiều cao khác nhau. Đáy của mọi bể đều nằm ở mặt đất.

Mỗi bể được nối với bể đứng trước nó trong dãy bằng một ống, và nối với bể
đứng sau nó bằng một ống khác. Các ống nối giữa các bể nằm ngang và có độ cao
tăng dần; tức là ống nối bể $i$ với bể $i+1$ nằm cao hơn ống nối bể $i$ với
bể $i-1$. Bể $1$ để hở, nên không khí và nước có thể tự do đi vào nó từ phía
trên. Tất cả các bể còn lại đều kín, nên không khí và nước chỉ có thể ra vào
qua các ống nối. Các ống nối đủ lớn để nước và không khí có thể chảy tự do và
đồng thời qua chúng, nhưng đủ nhỏ để kích thước của chúng được bỏ qua trong
bài toán này.

Dãy bể được đổ đầy bằng cách rót nước chậm vào miệng bể $1$, tiếp tục cho đến
khi mực nước đạt tới đỉnh của bể $1$. Khi mực nước dâng cao hơn các ống nối,
nước sẽ chảy giữa các bể.

Aqua Container Management cần một chương trình tính số mét khối nước có thể
đổ vào dãy bể trước khi mực nước đạt tới đỉnh của bể $1$.

Trong trường hợp đơn giản chỉ có hai bể như hình minh họa trong bản gốc, sau
khi quá trình đổ nước kết thúc, không khí ở phần trên của bể thứ hai bị nén
và áp suất của nó lớn hơn một atmosphere. Vì vậy mực nước trong bể thứ hai thấp
hơn mực nước trong bể thứ nhất.

\section*{Các nguyên lý vật lý}

Những nguyên lý vật lý sau đây hữu ích để giải bài toán này. Một số nguyên lý
là các xấp xỉ được chấp nhận cho mục đích của bài toán.

\begin{itemize}
  \item Nước chảy xuống chỗ thấp hơn.
  \item Trong một không gian hở, áp suất không khí bằng một atmosphere.
  \item Không khí có thể nén được: thể tích mà một lượng không khí cho trước
  chiếm phụ thuộc vào áp suất. Nước không nén được: thể tích mà một lượng nước
  cho trước chiếm là không đổi, không phụ thuộc vào áp suất.
  \item Áp suất không khí là như nhau tại mọi nơi trong một không gian kín.
  Nếu thể tích của không gian kín thay đổi, tích của thể tích và áp suất không
  khí trong không gian đó vẫn không đổi. Chẳng hạn, giả sử một khoang khí kín
  ban đầu có thể tích $V_1$ và áp suất $P_1$. Nếu thể tích của khoang khí đổi
  thành $V_2$, thì áp suất mới $P_2$ thỏa mãn
  \[
    P_1V_1 = P_2V_2.
  \]
  \item Trong một cột nước nằm dưới một khoang khí, áp suất nước tại độ sâu
  $D$ mét dưới mặt nước bằng áp suất không khí tại mặt nước cộng với
  $0.097D$ atmosphere. Điều này đúng bất kể khoang khí là hở hay kín.
  \item Trong một khối nước liên thông, ví dụ khi hai hay nhiều bể được nối
  với nhau bằng các ống nằm dưới mực nước, áp suất nước là hằng số tại mọi
  điểm có cùng độ cao.
\end{itemize}

\section*{Dữ liệu vào}

Dữ liệu vào gồm nhiều bộ kiểm thử, mỗi bộ biểu diễn một dãy bể nước khác
nhau. Mỗi bộ kiểm thử có ba dòng dữ liệu.

Dòng đầu tiên chứa một số nguyên $N$ $(2\le N\le 10)$, là số bể trong bộ kiểm
thử.

Dòng thứ hai chứa $N$ số thực dương, là chiều cao tính bằng mét của các bể
$1$ đến $N$. Dòng thứ ba chứa $N-1$ số thực. Trên dòng này, số thứ $k$ biểu
diễn độ cao so với mặt đất của ống nối bể $k$ và bể $k+1$. Các số trên dòng
thứ ba tăng dần, tức mỗi số lớn hơn số đứng trước nó.

Sau bộ kiểm thử cuối cùng là một dòng chứa số nguyên $0$.

\section*{Dữ liệu ra}

Với mỗi bộ kiểm thử, in một dòng chứa số thứ tự bộ kiểm thử, bắt đầu từ $1$,
theo sau là lượng nước tính bằng mét khối có thể đổ vào bể $1$ trước khi mực
nước đạt tới đỉnh của bể $1$. In kết quả với ba chữ số sau dấu thập phân.

In một dòng trống sau phần kết quả của mỗi bộ kiểm thử. Sử dụng đúng định
dạng của ví dụ.

\section*{Ví dụ}

\subsection*{Dữ liệu vào mẫu}

\begin{verbatim}
2
10.0 8.0
4.0
0
\end{verbatim}

\subsection*{Dữ liệu ra mẫu}

\begin{verbatim}
Case 1: 15.260
\end{verbatim}

\section*{Ghi chú nguồn}

Tokyo 2006--2007.

Tests-Setter: Derek Kisman.

\end{document}
